\section{Exercise 4 (Bonus): Music Transmission via USRP}
\label{sec:lab2_ex4_bonus}

This bonus exercise uses the same AM USRP transmitter and receiver chain from
Exercise~3, but replaces the single-tone message with a music waveform. The aim
is not to introduce a new analytical task. Instead, it serves as an end-to-end
verification that the hardware chain can carry a richer audio signal while
preserving recognizable structure after demodulation.

The received audio appears consistent with the melody of \emph{Fur Elise} by
Beethoven, which suggests that the modulation, transmission, and demodulation
stages remain functional for a more realistic baseband input.

\begin{figure}[H]
\centering
\labincludegraphics[width=0.8\linewidth]{figures/AM_Music_USRP.png}
\caption{USRP AM verification capture using a music waveform.}
\label{fig:lab2_ex4_music}
\end{figure}

\subsection{Interpretation}

From the plot alone, the main conclusion is that the recovered signal behaves
like a genuine audio waveform rather than a single sinusoid. This is enough for
verification of the transmission chain, but it is not strong evidence for
identifying the exact song from spectrum or PSD information alone. Melody is a
time-varying property, so exact identification would require the recovered audio
itself or a time-frequency representation such as a spectrogram.
